\section{Proofs for Fiber Information Geometry and Inference Histories}
\label{sec:proofs-fiber-geometry}

\noindent\emph{Proof of \Cref{prop:pb-statistical-tensor-descent}.}
Each $\widehat\rho_x(g)$ is pushforward by the bimeasurable sample map $\rho_x(g)$, which
does not depend on the law parameter.  Pushforward is a unitary map
$L^2(p)\to L^2(\rho_x(g)_\#p)$, $f\mapsto f\circ\rho_x(g)^{-1}$, differentiability in
quadratic mean makes $\mathbf s_Z\circ\rho_x(g)^{-1}$ the score of the pushed direction,
and the change of variables preserves \eqref{eq:pb-fiber-fisher} and
\eqref{eq:pb-fiber-amari}; since $\iota_{u\cdot g}=\iota_u\circ\widehat\rho_x(g)$,
$(\iota_u^{-1})^*\tau$ is independent of $u$ exactly when
$\widehat\rho_x(g)^*\tau=\tau$ for every $g$.  For the counterexample,
$\mathrm X=\mu/\sqrt2$ and $\mathrm S=\varsigma$ turn $g^F$ into
$2(d\mathrm X^2+d\mathrm S^2)/\mathrm S^2$, twice the hyperbolic metric of the upper
half-plane, of which the real M\"obius map $z\mapsto(1+z)/(1-z)$ of determinant $2>0$ is
an isometry.  It fixes $z=\mathrm i$ with derivative $2/(1-\mathrm i)^2=\mathrm i$, a
quarter turn, so it sends the $g^F$-unit vector $\partial_\mu$ at $(0,1)$ to
$\partial_\varsigma/\sqrt2$.  There the only nonzero components of $T^F$ are
$T^F(\partial_\mu,\partial_\mu,\partial_\varsigma)=2$ and
$T^F(\partial_\varsigma,\partial_\varsigma,\partial_\varsigma)=8$, so the cubic values of
the two unit vectors are $0$ and $8/2^{3/2}=2\sqrt2$.  $\square$

\noindent\emph{Proof of \Cref{prop:pb-kl-divergence-jets}.}
Invariance is \eqref{eq:cg-kl-bimeasurable-invariance} for the sample maps $\rho(g)$.
With local parameters $\lambda$ and $\lambda'$ of the two arguments, primed derivatives
acting on $\lambda'$, and $\mathbf s_i=\partial_i\log p_\lambda$, the identities
$\E_\lambda[\mathbf s_i]=0$, $\partial_ip_\lambda=p_\lambda\mathbf s_i$, and
$\partial_i\partial_jp_\lambda=p_\lambda(\mathbf s_i\mathbf s_j+\partial_i\mathbf s_j)$
and differentiation at $\lambda'=\lambda$ under \cref{hyp:pb-regular-models} give
\begin{align}
-\partial_i\partial_j'\KL(p_\lambda\Vert p_{\lambda'})
&=\E_\lambda[\mathbf s_i\mathbf s_j],\notag\\
-\partial_i\partial_j\partial_k'\KL(p_\lambda\Vert p_{\lambda'})
&=\E_\lambda[(\mathbf s_i\mathbf s_j+\partial_i\mathbf s_j)\mathbf s_k],
\notag\\
-\partial_i'\partial_j'\partial_k\KL(p_\lambda\Vert p_{\lambda'})
&=\E_\lambda[\mathbf s_k\partial_i\mathbf s_j].
\label{eq:pb-kl-jet-calculation}
\end{align}
The second derivatives of the score are symmetric, so the last two lines differ by
$\E_\lambda[\mathbf s_i\mathbf s_j\mathbf s_k]$, and positivity of the metric jet is
positivity of $g^F$.  The same dualistic structure is derived from contrast functions in
\citet{AmariNagaoka2000,AyJostLeSchwachhofer2017}.  $\square$

\noindent\emph{Proof of \Cref{cor:pb-orbit-fisher-form}.}
Positive definiteness of $g^F$ identifies the radical of $\mathfrak k_q$ with the kernel
of $\xi\mapsto\xi_{\mathcal B}(q)$, which is $\mathfrak g_q$.  The curve
$\gamma(t)=\widehat\rho(\exp t\xi)q$ is the integral curve of $\xi_{\mathcal B}$ through
$q$, so by uniqueness of integral curves $\xi_{\mathcal B}(q)=0$ exactly when
$\gamma\equiv q$, that is, when $\exp t\xi\in G_q$ for every $t$; this gives the
vanishing clause, and since $G_q$ is closed, its Lie algebra is the set of $\xi$ with
$\exp t\xi\in G_q$ for every $t$, which is $\mathfrak g_q$.  The
function $\phi(t)=\mathscr D(q,\gamma(t))$ is $C^3$ with $\phi(0)=\phi'(0)=0$ because
$\mathscr D(q,\cdot)$ is minimal at $q$, and differentiating
$\partial'_j\mathscr D(\lambda,\lambda)=0$ along the diagonal gives
$\partial'_i\partial'_j\mathscr D=-\partial_i\partial'_j\mathscr D=(g_{\mathscr D})_{ij}$
there, so $\phi''(0)=g_{\mathscr D}(\xi_{\mathcal B}(q),\xi_{\mathcal B}(q))$ and
Taylor's theorem gives the expansion; for relative entropy $g_{\mathscr D}=g^F$ by
\cref{prop:pb-kl-divergence-jets}.  $\square$

\noindent\emph{Proof of \Cref{thm:cg-fisher-contraction}.}
Pollard's Le~Cam decomposition at $\lambda_0$ is
$P_{\lambda_0+\delta}=\mathsf p_\delta P_{\lambda_0}+P_\delta^\perp$ with
$P_\delta^\perp\perp P_{\lambda_0}$, $P_\delta^\perp(\mathsf Y)=o(\lVert\delta\rVert^2)$,
and $\int(\sqrt{\mathsf p_\delta}-1-\tfrac12\delta^\top\mathbf s_{\lambda_0})^2
dP_{\lambda_0}=o(\lVert\delta\rVert^2)$.  Hence
$P_{\lambda_0+\delta}\otimes K=\mathsf p_\delta(y)(P_{\lambda_0}\otimes
K)+P_\delta^\perp\otimes K$, where the second term has mass $o(\lVert\delta\rVert^2)$ and
is singular, because a $P_{\lambda_0}$-null set carrying $P_\delta^\perp$, multiplied by
$\bar{\mathsf Y}$, carries it and is $P_{\lambda_0}\otimes K$-null since each
$K(y,\cdot)$ is a probability, and the square-root remainder of the first term is
unchanged because the kernel integrates to one.  The joint experiment is therefore
differentiable in quadratic mean with score $\mathbf s_{\lambda_0}(Y)$, without a common
dominating measure, and Pollard's preservation theorem for measurable statistics
\citep[Theorem~7]{Pollard2013}, applied to $(y,\bar y)\mapsto\bar y$, gives
differentiability of $\{P_\lambda K\}$ with the score \eqref{eq:cg-score-projection} in
$L^2(P_{\lambda_0}K;\R^d)$; the projected singular term may acquire an absolutely
continuous part, of mass still $o(\lVert\delta\rVert^2)$.  The law of total covariance
gives \eqref{eq:cg-fisher-loss}, and since conditional expectation is an orthogonal
projection, equality in the direction $z$ holds exactly when
$z^\top\mathbf s_{\lambda_0}(Y)$ is $\bar Y$-measurable.  $\square$

\noindent\emph{Proof of \Cref{thm:hist-retained-fisher-quotient}.}
Since $h_Q(Z,\cdot)=g(d\operatorname{ret}_QZ,d\operatorname{ret}_Q\cdot)$, the
Cauchy--Schwarz inequality for the semidefinite $g$ makes $h_Q(Z,Z)=0$,
$h_Q(Z,\cdot)=0$, and $d\operatorname{ret}_QZ\in\operatorname{rad}g$ equivalent, which is
\eqref{eq:hist-cds-fisher-radical} and the transversality criterion.  The radical is the
kernel of the constant-rank $C^1$ bundle map $h^\flat:T\mathcal Q\to T^*\mathcal Q$,
hence a $C^1$ subbundle, and $h$ descends to the quotient because it annihilates its
radical.  When $\operatorname{rad}h=\ker d\operatorname{ret}$, $\operatorname{ret}$ has
zero derivative along, hence is constant on, each connected leaf and factors as $\bar r\circ\pi$ through the leaf space,
with $\pi$ a submersion, $\ker d\pi=\ker d\operatorname{ret}$, and
$\bar r=\operatorname{ret}\circ\iota$ smooth for smooth local sections $\iota$ of $\pi$;
then $h=\pi^*(\bar r^*g)$, $d\bar r$ is injective, and
$\operatorname{im}d\bar r=\operatorname{im}d\operatorname{ret}$ meets
$\operatorname{rad}g$ only in zero, so $\bar r^*g$ is positive definite.  For a finite
family, $h=\sum_y\partial p_y\partial p_y/p_y$ is the Gram matrix of the Jacobian of
$\operatorname{ret}$ weighted by $1/p_y$, hence $C^2$ and positive definite exactly
where the Jacobian has full column rank; an immersion is locally injective because the
inverse function theorem \citep[Ch.~4]{Lee2012} makes a coordinate projection of
$\operatorname{ret}$ with invertible square Jacobian block a local diffeomorphism, and
the Levi-Civita connection of a $C^2$ metric is $C^1$.  For softmax,
$Z^\top h(z)Z=\operatorname{Var}_p(Z)$ vanishes exactly on constant vectors, so
$\operatorname{rad}h=\R\mathbf 1=\ker d\operatorname{ret}$ with constant rank $k-1$, and
an objective of $p$ has $d\Fenergy(\mathbf 1)=0$, so
$\nabla\Fenergy\in\mathbf 1^\perp=\operatorname{im}h$, which gives the solution class and
$h^+\nabla\Fenergy\in\operatorname{im}h=\mathbf 1^\perp$.  $\square$

\noindent\emph{Proof of \Cref{prop:hist-general-score-residual}.}
The tower property gives $\E[\mathbf s\mathbf s_i^\top]=\mathsf G_i$ and its transpose,
so expanding $\mathfrak K$ gives $\mathfrak K=\mathsf G_{\mathrm J}-\sum_i\mathsf
G_i+\mathfrak C$, which rearranges to \eqref{eq:hist-joint-weighted-marginal-residual};
the criterion is its restriction, and \eqref{eq:hist-joint-versus-product} is its
quadratic form for $n=2$ and unit weights.  Under independence the joint score is the
sum of the block scores, so $\mathbf r=0$, $\mathfrak C=0$, and the weight condition is
$(1-w_1)\mathsf G_1+(1-w_2)\mathsf G_2=0$ in the stated direction.  For the normal
location family with precision $\Lambda$ the joint Fisher metric is $\Lambda$ and each
marginal is a location family of variance $1/(1-\mathrm c^2)$, so
$\mathsf G_1+\mathsf G_2=(1-\mathrm c^2)\mathrm I_2$, which gives the displayed difference
and its two values.  The lift metrics are the pullbacks of $d\lambda_1^2+d\lambda_2^2$
along the three maps.  $\square$

\noindent\emph{Proof of \Cref{thm:hist-three-agent-parity-lift}.}
Each atom has $p_0/\mathrm D=\prod_{y_{i,x}=1}(1-\lambda_{i,x})^{-1}
\prod_{y_{i,x}=0}\lambda_{i,x}^{-1}>1$, so
$p_\varkappa\geq p_0-|\varkappa|\mathrm D>(1-|\varkappa|)\mathrm D>0$.  Summing out any
one coordinate cancels the parity term, because $\mathrm D$ does not depend on $y$ and
$\operatorname{par}$ changes sign with that coordinate; this gives normalization and the
product proper marginals, hence $\operatorname{mar}\circ\iota_\varkappa=\operatorname{id}$,
which with polynomial dependence gives smoothness and injectivity of
$d\iota_\varkappa$.  The marginal map is the restriction to the open $63$-dimensional
simplex of a linear map of rank six, so its fibers are open subsets of $57$-dimensional
affine spaces, and $p_\varkappa$ is affine in $\varkappa$.  A paired complement changes
two bits, so it preserves $\operatorname{par}$ and $\mathrm D$ while $p_0$ transforms by
Bernoulli relabeling, and a typed permutation preserves all three expressions, which
proves \eqref{eq:hist-parity-paired-equivariance}; a single complement changes the sign
of $\operatorname{par}$ and sends $\lambda_{i,x}$ to $1-\lambda_{i,x}$, which preserves
$\mathrm D$ and proves \eqref{eq:hist-parity-odd-pseudoscalar}.  $\square$

\noindent\emph{Proof of \Cref{cor:hist-parity-fisher-residual}.}
Two-bit marginals are products by \cref{thm:hist-three-agent-parity-lift}, so centered
scores of distinct bits are uncorrelated.  With $\mathbf s_{i,x}$ the one-bit Bernoulli
score, differentiation gives the residual components
\begin{equation}
\mathbf r_{i,x}
=\frac{\varkappa\operatorname{par}(y)\mathrm D(\lambda)}{p_\varkappa(y;\lambda)}
\left(\frac{1-y_{i,x}}{\lambda_{i,x}}-\frac{y_{i,x}}{1-\lambda_{i,x}}\right).
\label{eq:hist-parity-score-residual-coordinate}
\end{equation}
If $Z^\top\mathbf r$ vanished on every atom, positivity and $\varkappa\mathrm D\ne0$
would make the $Z$-weighted sum of the six bracketed one-bit functions vanish on the
cube, and comparing two atoms that differ in one coordinate forces each component of $Z$
to vanish; so $\mathfrak K\succ0$.  At the center $\mathrm D=4^{-6}$, its derivative
vanishes, and the law is \eqref{eq:hist-parity-center-law},
so, each one-bit score having square $4$, averaging
$(1+\bar\varkappa\operatorname{par})^{-1}=(1-\bar\varkappa\operatorname{par})/(1-\bar\varkappa^2)$
over the cube gives \eqref{eq:hist-parity-center-fisher}, and $\bar\varkappa=1/128$ gives
\eqref{eq:hist-parity-half-fisher-residual}.  The one-bit tensors at the center are
$4e_{i,x}e_{i,x}^\top$, which gives \eqref{eq:hist-parity-center-weight-equality}, and the
last sentence compares $\varkappa=0$ with $\varkappa\ne0$.  $\square$

\noindent\emph{Proof of \Cref{prop:hist-parity-hyperedge-vfe}.}
Half of the atoms have each parity, so averaging \eqref{eq:hist-parity-hyperedge-kernel}
under $\operatorname{Unif}$ gives evidence $\tfrac12$, Bayes' rule gives the posterior
$64^{-1}(1+\bar\varkappa\operatorname{par})$, which is \eqref{eq:hist-parity-center-law},
and the gap identity of the free energy gives the tight value $-\log\tfrac12$.  Every
posterior one-bit marginal is uniform, so its product law is $\operatorname{Unif}$, and
since the two parities have equal mass,
$-\E_{\operatorname{Unif}}\log K_{\mathrm{hyp}}(1\given Y)
=\log2-\tfrac12\log(1-\bar\varkappa^2)$, which is
\eqref{eq:hist-parity-product-lift-excess}.  The posterior and $\operatorname{Unif}$ have
the same one-bit marginals, the posterior input and the evidence are unchanged, and the
two free-energy values differ when $\varkappa\ne0$.  $\square$

\noindent\emph{Proof of \Cref{thm:hist-fisher-clock-invariance}.}
The chain rule gives $\dot{\widetilde Q}=(\dot Q\circ\phi)\phi'$ with $\phi'>0$, and the
substitution $t=\phi(\widetilde t)$ proves length invariance and
\eqref{eq:hist-clock-reparameterization-law}.  Positive speed makes the clock invertible,
and dividing by $d\mathfrak t/dt$ gives unit speed.  Substituting
\eqref{eq:hist-vfe-ray-representative} into \eqref{eq:hist-fisher-speed} gives
\eqref{eq:hist-vfe-clock-rate}, and the other two identities follow from
\eqref{eq:hist-fisher-gradient-definition}.  For a second metric,
$d\Fenergy(Z)=\mathsf G'(\operatorname{grad}^{\mathsf G'}\Fenergy,Z)$, and the
Cauchy--Schwarz inequality gives \eqref{eq:hist-clock-metric-dissipation-bound}.  For a
block-diagonal metric the squared speeds add, so
$\dot{\mathfrak t}_i\leq\dot{\mathfrak t}\leq\sum_i\dot{\mathfrak t}_i$, and Minkowski's
integral inequality for the vector $(\dot{\mathfrak t}_i)_i$ gives
$\bigl(\sum_i(\int\dot{\mathfrak t}_i)^2\bigr)^{1/2}\leq\int\dot{\mathfrak t}$.  $\square$

\noindent\emph{Proof of \Cref{prop:hist-analytic-finite-duration}.}
Along the history $\tfrac{d}{dt}\Fenergy(\gamma)=-\lVert\dot\gamma\rVert_{\mathsf
G}^2\leq0$, so $\Fenergy(\gamma(t))$ decreases to $\Fenergy_*=\Fenergy(Q_*)$; if it
reaches $\Fenergy_*$ at some $t_1$, then $\dot\gamma=0$ on $[t_1,\infty)$ and the
duration is that of a compact segment.  Otherwise choose a closed Euclidean ball $\bar B$
of radius $\delta_0$ around $Q_*$ in the chart on which
$\mathrm c_1\mathrm I\preceq\mathsf G\preceq\mathrm c_2\mathrm I$ with
$0<\mathrm c_1\leq\mathrm c_2$ and the {\L}ojasiewicz gradient inequality
$\lVert\nabla\Fenergy\rVert\geq\mathrm c_0|\Fenergy-\Fenergy_*|^{\mathrm q}$ holds with
$\mathrm c_0>0$ and $\mathrm q\in[0,1)$ \citep[Lemma~2.1]{absil2005convergence}.  While
$\gamma$ stays in $\bar B$,
$\lVert\dot\gamma\rVert_{\mathsf G}^2=\nabla\Fenergy^\top\mathsf G^{-1}\nabla\Fenergy
\geq\mathrm c_2^{-1}\lVert\nabla\Fenergy\rVert^2$, hence
\begin{equation}
\frac{d}{dt}\left(\Fenergy(\gamma)-\Fenergy_*\right)^{1-\mathrm q}
=-(1-\mathrm q)\left(\Fenergy(\gamma)-\Fenergy_*\right)^{-\mathrm q}
\lVert\dot\gamma\rVert_{\mathsf G}^2
\leq-(1-\mathrm q)\mathrm c_0\mathrm c_2^{-1/2}\lVert\dot\gamma\rVert_{\mathsf G},
\label{eq:hist-lojasiewicz-length-estimate}
\end{equation}
so, with $\mathrm C=((1-\mathrm q)\mathrm c_0\mathrm c_2^{-1/2})^{-1}$, a segment
$\gamma([t_1,t_2])\subseteq\bar B$ has Fisher length at most
$\mathrm C(\Fenergy(\gamma(t_1))-\Fenergy_*)^{1-\mathrm q}$ and Euclidean length at most
$\mathrm c_1^{-1/2}$ times that.  For $t_1$ with $\lVert\gamma(t_1)-Q_*\rVert<\delta_0/2$
and $\mathrm c_1^{-1/2}\mathrm C(\Fenergy(\gamma(t_1))-\Fenergy_*)^{1-\mathrm
q}<\delta_0/2$, a first exit from $\bar B$ after $t_1$ would give a segment in $\bar B$ of
Euclidean length both at least and less than $\delta_0/2$; so $\gamma$ stays in $\bar B$,
has finite length on $[t_1,\infty)$, and is Cauchy, hence converges to $Q_*$.  If $Q_*$
were noncritical, $\tfrac{d}{dt}\Fenergy(\gamma)$ would tend to
$-\lVert\operatorname{grad}^{\mathsf G}\Fenergy(Q_*)\rVert^2_{\mathsf G}<0$,
contradicting convergence of $\Fenergy(\gamma)$.  This is the argument of
\citet[Theorem~2.2]{absil2005convergence}, with its angle condition supplied by the
metric bounds.  In the smooth counterexample the Euclidean plane is the Fisher geometry
of the unit-variance normal location family, and a curve of finite length in the plane
converges, so a history that does not converge has infinite duration.  $\square$

\noindent\emph{Proof of \Cref{prop:hist-invariant-history-horizontal}.}
Differentiating $\Fenergy(\exp(t\xi)\cdot Q)=\Fenergy(Q)$ at $t=0$ gives
$\mathsf G(\operatorname{grad}^{\mathsf G}\Fenergy,\xi_{\mathcal Q}(Q))=0$, and for a
velocity $Z\perp V_Q$ the quantity $\lVert Z-\xi_{\mathcal Q}(Q)\rVert^2=\lVert
Z\rVert^2+\lVert\xi_{\mathcal Q}(Q)\rVert^2$ is minimal where $\xi_{\mathcal Q}(Q)=0$.
In the free and proper case the orbit space is a smooth manifold and $\pi$ a surjective
submersion \citep[Theorem~21.10]{Lee2012}.  With
$\operatorname{Hor}_Q=(\ker T_Q\pi)^\perp$, $T_Q\pi$ maps $\operatorname{Hor}_Q$
isomorphically onto $T_{\pi(Q)}(\mathcal Q/\mathbb G)$; define $\bar{\mathsf G}$ by
transporting $\mathsf G|_{\operatorname{Hor}_Q}$ through this isomorphism.  Each
$L_g:Q\mapsto g\cdot Q$ is an isometry with $\pi\circ L_g=\pi$, so $TL_g$ maps
$\operatorname{Hor}_Q$ onto $\operatorname{Hor}_{gQ}$ with $T_{gQ}\pi\circ TL_g=T_Q\pi$
and the transported inner products agree along orbits; $\bar{\mathsf G}$ is smooth
because, through a smooth local section $\iota_\pi$ of $\pi$ given by the rank theorem,
its coefficients are those of $\mathsf G$ on the horizontal projections of the pushed
coordinate fields, and it is the unique metric for which $\pi$ is a Riemannian
submersion.  The invariant $\Fenergy$ is constant on the fibers of $\pi$, so
$\bar\Fenergy=\Fenergy\circ\iota_\pi$ is well defined and $C^1$ with
$\Fenergy=\bar\Fenergy\circ\pi$; hence $d\Fenergy$ vanishes on $\ker T\pi$, the gradient
is horizontal, and for horizontal $Z$,
$\mathsf G(\operatorname{grad}^{\mathsf G}\Fenergy,Z)=d\bar\Fenergy(T\pi Z)
=\bar{\mathsf G}(\operatorname{grad}^{\bar{\mathsf G}}\bar\Fenergy,T\pi Z)$, which gives
the gradient identity.  Applying $T\pi$ to \eqref{eq:hist-vfe-ray-representative} shows
that $\pi\circ Q$ is a natural-gradient history of $\bar\Fenergy$, and since $T\pi$ is
isometric on horizontal vectors and the velocity is horizontal, the two Fisher speeds
agree.  $\square$

\noindent\emph{Proof of \Cref{thm:hist-global-clock-exactness}.}
Exact forms are closed and have zero periods; conversely, under
\eqref{eq:hist-global-clock-period-criterion}, integrating $\bm\alpha$ from a base point
along curves in $\mathcal U$ is path independent and defines a function with
$d\mathfrak t^\perp=\bm\alpha$.  Differentiating $\bm\alpha=-\mathsf n^{-1}d\Fenergy$
gives \eqref{eq:hist-normalized-form-curvature}.  Since $d\Fenergy\ne0$, the form
$d\mathsf n\wedge d\Fenergy$ vanishes exactly when $d\mathsf n$ is locally a multiple of
$d\Fenergy$, that is, when $\mathsf n$ does not depend on the other coordinates of a chart
whose first coordinate is $\Fenergy$.  Then $d\mathfrak t^\perp\wedge
d\Fenergy=-\mathsf n^{-1}d\Fenergy\wedge d\Fenergy=0$, so $\mathfrak t^\perp$ is constant
on connected pieces of level sets and locally $\mathfrak t^\perp=\tau\circ\Fenergy$ with
$\tau'=-1/\mathsf n<0$.  Since $\bm\alpha(\mathbf e_{\Fenergy})=\mathsf G(\mathbf
e_{\Fenergy},\mathbf e_{\Fenergy})=1$, the unit rate follows, and the kernel of
$\bm\alpha=\mathbf e_{\Fenergy}^\flat$ is the orthogonal complement of
$\mathbf e_{\Fenergy}$.  In the example $\mathsf n=(z_1^2+z_2^2)^{1/2}$ and
$\bm\alpha=-(z_2dz_1+z_1dz_2)/\mathsf n$, whose exterior derivative is
\eqref{eq:hist-nonexact-clock-example}, with coefficient vanishing only on the diagonal.
$\square$

\noindent\emph{Proof of \Cref{prop:hist-cross-section-clock}.}
The unit field $\mathbf e_{\Fenergy}$ is $C^1$ because $\mathsf n>0$.  Let
$\mathcal W\subseteq\R\times\mathcal S$ be the open set of pairs $(t,q)$ for which
$\operatorname{Fl}^{\mathbf e_{\Fenergy}}_{t}(q)$ is defined in $\mathcal U$, and put
$\Xi(t,q)=\operatorname{Fl}^{\mathbf e_{\Fenergy}}_{t}(q)$, a $C^1$ map.  Each orbit meets
$\mathcal S$, so $\Xi$ is onto; $\Fenergy$ decreases strictly along orbits, so no orbit
is closed, and each orbit meets $\mathcal S$ once, so $\Xi$ is injective.  At $(t,q)$,
$T\Xi$ sends $\partial_{t}$ to $\mathbf e_{\Fenergy}$ and $T_q\mathcal S$ onto the image
of a complement of $\mathbf e_{\Fenergy}(q)$ under the invertible map
$T\operatorname{Fl}^{\mathbf e_{\Fenergy}}_{t}$, which also carries
$\mathbf e_{\Fenergy}(q)$ to $\mathbf e_{\Fenergy}$ at the image point, so transversality
makes $T\Xi$ invertible.  Hence $\Xi$ is a $C^1$ diffeomorphism, and
$\mathfrak t^{\mathcal S}=\operatorname{pr}_1\circ\Xi^{-1}$ satisfies
$d\mathfrak t^{\mathcal S}(\mathbf e_{\Fenergy})=1$, with $t$ the Fisher arc length
because $\mathbf e_{\Fenergy}$ has unit speed.  A level set is transversal and met once
because $d\Fenergy(\mathbf e_{\Fenergy})=-\mathsf n<0$.  If the level sets of
$\mathfrak t^{\mathcal S}$ are orthogonal to $\mathbf e_{\Fenergy}$, then
$d\mathfrak t^{\mathcal S}$ vanishes on the orthogonal complement and equals one on
$\mathbf e_{\Fenergy}$, so it equals $\bm\alpha$.  In the example the Euclidean descent
orbits $\dot z_1=-z_2$, $\dot z_2=-z_1$ in the quadrant are
$z_1=\mathrm c_1e^t+\mathrm c_2e^{-t}$, $z_2=-\mathrm c_1e^t+\mathrm c_2e^{-t}$ with
$\mathrm c_2>|\mathrm c_1|$, along which $\Fenergy=\mathrm c_2^2e^{-2t}-\mathrm
c_1^2e^{2t}$ decreases from $+\infty$ toward $0$, so each meets $\{z_1z_2=1\}$ exactly
once.  $\square$

\noindent\emph{Proof of \Cref{thm:hist-record-clock-contraction}.}
At each $t$, \cref{thm:cg-fisher-contraction} applied to the one-parameter family
$\{P_{t'}\}$ gives \eqref{eq:hist-record-score-projection} and the defect identity
\eqref{eq:hist-record-conditional-variance-defect} with its equality clause; taking
nonnegative square roots and integrating gives \eqref{eq:hist-record-clock-contraction}.
For a dominated channel, differentiating the joint density
$p_t(y)\mathsf k_t(\bar y\given y)$ gives the joint score
$\mathbf s_t(y)+\partial_t\log\mathsf k_t(\bar y\given y)$, whose projection onto the
record is \eqref{eq:hist-parameter-dependent-channel-score}.  In the example the fine
law does not move, and a Bernoulli law of success probability
$\mathrm b(t)=(1+e^{-t})^{-1}$ has Fisher speed
$|\dot{\mathrm b}|/\sqrt{\mathrm b(1-\mathrm b)}=\sqrt{\mathrm b(1-\mathrm b)}=(2\cosh(t/2))^{-1}$.
$\square$

\noindent\emph{Proof of \Cref{thm:hist-oriented-semiconjugacy}.}
Put $\gamma(t)=\mathsf F_\ell(\operatorname{Fl}^{\mathscr V_\ell}_tQ)$, so
$\dot\gamma=\mathsf a_\ell(\operatorname{Fl}^{\mathscr V_\ell}_tQ)\mathscr
V_{\ell+1}(\gamma)$.  The integrand of $\psi_Q$ is continuous and positive, so $\psi_Q$ is
a $C^1$ diffeomorphism onto an open interval.  The curve $\gamma\circ\psi_Q^{-1}$
satisfies $(\gamma\circ\psi_Q^{-1})'=\mathscr V_{\ell+1}(\gamma\circ\psi_Q^{-1})$ on
$\mathcal T'_Q$ and starts at $\mathsf F_\ell Q$, so by local Lipschitz continuity it is
contained in the maximal integral curve through $\mathsf F_\ell Q$; this gives the
inclusion and the flow identity, with the curve constructed before the coarse flow is
evaluated at $\psi_Q(t)$.  The endpoint formula is monotonicity of $\psi_Q$, and in the
complete case the limits are the two integrals.  In the examples
$\mathsf a_\ell=\mathsf F_\ell'$ because both fields are $\partial_z$.  The converse is
the chain rule.  $\square$

\noindent\emph{Proof of \Cref{prop:hist-natural-gradient-sufficiency}.}
In the first example, \eqref{eq:hist-oriented-vector-field-semiconjugacy} reads
$(z_1,2z_2)=\mathsf a(z_1,2z_2/\mathrm c)$, which forces $\mathsf a=1$ and $\mathrm c=1$
where $z_1z_2\ne0$.  Differentiating the compatibility identity gives
$d\Fenergy_\ell(Z)=\chi_\ell'd\Fenergy_{\ell+1}(T\mathsf F_\ell Z)=0$ for
$Z\in\ker T\mathsf F_\ell$, so $\Fenergy_\ell$ is constant along curves tangent to the
fibers, $Z_0=\operatorname{grad}^{\mathsf G_\ell}\Fenergy_\ell$ lies in
$\operatorname{Hor}_Q$, and for $Z\in\operatorname{Hor}_Q$
\[
\mathsf G_{\ell+1}(T\mathsf F_\ell Z_0,T\mathsf F_\ell Z)
=\mathsf d_\ell^2d\Fenergy_\ell[Z]
=\mathsf d_\ell^2\chi_\ell'\mathsf G_{\ell+1}\left(\operatorname{grad}^{\mathsf
G_{\ell+1}}\Fenergy_{\ell+1}\circ\mathsf F_\ell,T\mathsf F_\ell Z\right).
\]
Since $T\mathsf F_\ell(\operatorname{Hor}_Q)$ is the whole tangent space,
$T\mathsf F_\ell Z_0=\mathsf d_\ell^2\chi_\ell'\operatorname{grad}^{\mathsf
G_{\ell+1}}\Fenergy_{\ell+1}\circ\mathsf F_\ell$, and negating gives the claim.  The
quotient case is $\mathsf d_\ell\equiv1$ and $\chi_\ell=\operatorname{id}$.  In the
normal example the kernel sends $\mathcal N(\mu,1)$ to $\mathcal N(\mu,2)$, the two
Fisher metrics are $d\mu^2$ and $\tfrac12d\mu^2$, and $\mathscr V_\ell=-\mu$,
$\mathscr V_{\ell+1}=-2\mu$.  For the exact-contraction objective,
$\Fenergy_{\ell+1}\circ\mathsf F_\ell\leq\Fenergy_\ell$ with equality exactly on the set
$\mathcal A$ of fiberwise minimizers, so compatibility with
$\chi_\ell=\operatorname{id}$ holds exactly on $\mathcal A$; an open subset of
$\mathcal A$ would contain, with each of its points, a neighborhood of that point in its
fiber, which has more than one point, so $\mathcal A$ has empty interior.  In the last
example $\mathcal A=\{f=s\}$ and $d\Fenergy_\ell(\partial_f)=f-s$, which vanishes on no
open set, so the fiber condition fails for every $\chi_\ell$.  $\square$
